To evaluate the efficiency gain from heterogeneous parallel simulation, we conduct controlled experiments under five representative settings: RoboDojo with heterogeneous or non-heterogeneous parallel simulation under zero-action rollouts, RoboDojo with heterogeneous or non-heterogeneous parallel simulation coupled with $\pi_{0.5}$ policy inference, and RoboTwin 2.0 with zero-action rollouts.
All simulation and rendering workloads are executed on 8$\times$RTX 4090 GPUs.
In our setup, each RTX 4090 runs one RoboDojo simulation process with 10 parallel environments, while RoboTwin 2.0 runs two simulation processes per GPU.
Notably, RoboTwin 2.0 renders three camera views at a resolution of $320 \times 240$, whereas RoboDojo renders at $640 \times 480$, resulting in a higher per-frame rendering workload.
For non-zero-action RoboDojo settings, $\pi_{0.5}$ inference is performed on an A800 server within the same local network.
The zero-action setting measures raw simulation throughput without policy inference overhead, whereas the $\pi_{0.5}$ setting evaluates end-to-end efficiency under a realistic remote policy-deployment scenario.

As shown in Table~\ref{tab:simulation_efficiency}, heterogeneous parallel simulation substantially improves evaluation throughput. 
Under zero-action rollouts, RoboDojo achieves 77.4 interactions/s, compared with 40.0 interactions/s under non-heterogeneous parallel simulation, yielding a $1.94\times$ speedup. 
When $\pi_{0.5}$ inference is included, heterogeneous parallel simulation achieves 64.0 interactions/s, outperforming the non-heterogeneous setting by $1.63\times$. 
These results show that the proposed infrastructure improves both raw simulation throughput and practical end-to-end evaluation efficiency. 
Compared with RoboTwin 2.0 zero-action rollouts, RoboDojo achieves higher normalized throughput while supporting more diverse tasks and scenes. 
This efficiency enables large-scale policy evaluation and rapid feedback across comprehensive capability dimensions.

\begin{table}[h]
\centering
\caption{\textbf{Simulation evaluation efficiency.}
All simulation and rendering workloads are executed on 8$\times$RTX 4090 GPUs.
RoboDojo runs one simulation process with 10 parallel environments per RTX 4090, while RoboTwin 2.0 runs two simulation processes per GPU.
For non-zero-action RoboDojo settings, $\pi_{0.5}$ policy inference is performed on an A800 server within the same local network.}
\label{tab:simulation_efficiency}
\resizebox{0.9\textwidth}{!}{
\begin{tabular}{l p{0.45\textwidth} c c c}
\toprule
\textbf{Benchmark} & \textbf{Evaluation Setting} & \textbf{Frames} & \textbf{Time} & \textbf{Avg. Speed} \\
\midrule
RoboDojo 
& Heterogeneous parallel simulation, zero action 
& 1,640,000 & 5h 53m & 77.4 interactions/s \\

RoboDojo 
& Non-heterogeneous parallel simulation, zero action 
& 1,640,000 & 11h 22m & 40.0 interactions/s \\

RoboDojo 
& Heterogeneous parallel simulation + $\pi_{0.5}$ inference 
& 1,640,000 & 7h 07m & 64.0 interactions/s \\

RoboDojo 
& Non-heterogeneous parallel simulation + $\pi_{0.5}$ inference 
& 1,640,000 & 11h 38m & 39.2 interactions/s \\

RoboTwin 2.0 
& Zero-action rollout 
& 3,100,000 & 19h 19m & 44.6 interactions/s \\
\bottomrule
\end{tabular}
}
\end{table}